%% Angles remarquables de 0 à pi sur le cercle trigonométrique (manuel p. 98) :
%% radians (en noir) et degrés (en gris) — c'est le tableau de conversion sous forme visuelle.
%% Rayons du premier quadrant en rose, ceux du deuxième en gris.
%% Les étiquettes de pi/4 et 3pi/4 sont repoussées sur un rayon plus grand pour éviter
%% tout chevauchement avec leurs voisines.
\documentclass[tikz,border=4pt]{standalone}
\usepackage{liaisons}
\begin{document}
\begin{tikzpicture}[x=2cm,y=2cm,
  axe/.style={-{Stealth[length=5pt]}, line width=0.6pt},
  etiq/.style={align=center, font=\small, inner sep=2pt}]

%% ---------------- repère et cercle ------------------------------------
\draw[axe] (-1.13,0) -- (1.13,0);
\draw[axe] (0,-1.13) -- (0,1.13);
\draw[line width=1pt] (0,0) circle (1);
\node[font=\small, anchor=north east, inner sep=2pt] at (-0.03,-0.03) {$O$};

%% ---------------- rayons : premier quadrant (rose), deuxième (gris) ---
\foreach \angl in {30,45,60} {
  \draw[solideC, line width=0.9pt] (0,0) -- (\angl:1);
  \fill[solideC] (\angl:1) circle (1.6pt);}
\foreach \angl in {120,135,150} {
  \draw[black!55, line width=0.9pt] (0,0) -- (\angl:1);
  \fill[black!55] (\angl:1) circle (1.6pt);}
\foreach \angl in {0,90,180} { \fill (\angl:1) circle (1.6pt); }

%% ---------------- étiquettes : radians + degrés ------------------------
%% \angl / rayon de l'étiquette / ancrage / mesure en radians / mesure en degrés
\foreach \angl/\ray/\anc/\labr/\dgr in {%
    0/1.16/west/{0}/{0},
    30/1.16/west/{\frac{\pi}{6}}/{30},
    45/1.62/west/{\frac{\pi}{4}}/{45},
    60/1.16/west/{\frac{\pi}{3}}/{60},
    90/1.18/south/{\frac{\pi}{2}}/{90},
    120/1.16/east/{\frac{2\pi}{3}}/{120},
    135/1.62/east/{\frac{3\pi}{4}}/{135},
    150/1.16/east/{\frac{5\pi}{6}}/{150},
    180/1.16/east/{\pi}/{180}} {
  \node[etiq, anchor=\anc] at (\angl:\ray)
       {$\labr$ rad\\[-1pt] {\scriptsize\textcolor{black!55}{$\dgr^{\circ}$}}};}

%% ---------------- points I et J ---------------------------------------
\node[font=\small, anchor=north east, inner sep=3pt] at (1,0) {$I$};
\node[font=\small, anchor=north east, inner sep=3pt] at (0,1) {$J$};
\end{tikzpicture}
\end{document}
